\documentclass[10pt,a4paper]{amsart}
\usepackage[margin=19mm]{geometry}
\usepackage{amsmath,amssymb,mathtools}
\usepackage[scr=rsfs,cal=euler]{mathalfa}
\usepackage{xfrac}
\usepackage[unicode,hidelinks,bookmarks=false]{hyperref}
\newcommand{\ie}{i.e.\ }
\newcommand{\cf}{cf.\ }
\newcommand{\resp}{respectively\ }
\newcommand{\Subsection}{\subsection}
\providecommand{\nobreakdash}{\nobreak\hbox{-}}
\setlength{\emergencystretch}{2em}
\multlinegap0pt
\usepackage{bm}
\usepackage{bbm}
\usepackage{dsfont}
\usepackage{xspace}
\usepackage{cancel}
\usepackage{mathtools}
\usepackage{amsthm}
\makeatletter
\@ifundefined{theorem}{\newtheorem{theorem}{Theorem}}{}
\@ifundefined{proposition}{\newtheorem{proposition}[theorem]{Proposition}}{}
\makeatother
\title[Lacunary recurrences and 2-adic properties of Eisenstein ser]{Lacunary recurrences and 2-adic properties of Eisenstein series}
\author{Liubomir Chiriac, Andrei Jorza}
\date{Source: arXiv:2605.09738v2 (CC BY 4.0)}
\begin{document}
\maketitle

\noindent\emph{Source and licence.} Lacunary recurrences and 2-adic properties of Eisenstein series. Version arXiv:2605.09738v2 (CC BY 4.0), distributed under the Creative Commons Attribution 4.0 International licence, as recorded in the arXiv licence metadata for this exact version and shown on the abstract page https://arxiv.org/abs/2605.09738v2. Full rights notice: licenses/arxiv-2605.09738-CC-BY-4.0.txt.\par\smallskip

\noindent\textit{Editorial scope.} Counted unit: the Proposition whose source label is \texttt{Min-Witness} in the article (source0.tex lines 341--349, source label \texttt{Min-Witness}), delivered together with the article's own complete original proof of it (source0.tex lines 351--395, one \texttt{proof} environment of 45 lines). No mathematics is written, repaired or re-derived: statement and proof are the article's own text line for line. Required context retained uncounted: none. Every definition, display and label of those lines is retained unchanged. Omitted from this package and not redistributed: 1--340, 350--350, 396--607. Nothing inside a retained excerpt is omitted or altered: every retained body is delivered exactly as the article's lines are written, so no omission receipt of any kind accompanies this package. Citations: the retained excerpts carry no citation command at all, so no bibliography entry of the article is transcribed and no reference is redistributed. Reference adaptation: none was needed and none was made; every reference inside the delivered unit resolves inside it, so no unresolved reference or citation is produced. Printed slips kept literal: no printed word, symbol, hypothesis or number of the counted statement or of its proof was altered. Layout and class: the article's own class is \texttt{amsart} with packages including fullpage, amsmath, amssymb, amsthm, hyperref; the wrapper is the lane's pinned \texttt{amsart} editorial wrapper at 10pt on A4, which restates the theorem-like environments the retained text needs and adds the article's own macro definitions that the retained text needs (none), with extra wrapper packages (bm, bbm, dsfont, xspace, cancel, mathtools). The delivered numbering is the unit's own. Assets: no retained span contains a figure environment, a graphics-inclusion command, a PDF-page inclusion, a table or a TikZ picture; no image file is copied into this package, the package declares no asset dependency and no image file is redistributed with it. Licence: version arXiv:2605.09738v2 of this article is distributed under the Creative Commons Attribution 4.0 International licence, as recorded in the arXiv licence metadata for this exact version and shown on the abstract page https://arxiv.org/abs/2605.09738v2. A full rights notice is delivered at licenses/arxiv-2605.09738-CC-BY-4.0.txt. Characters: every retained excerpt is pure ASCII, so the delivered engine, XeLaTeX with its default Latin Modern text font, reports no missing character.
\begin{proposition} \label{Min-Witness}
Let $k\ge 6$ be even and not a power of 2. 

	\begin{enumerate}
		\item[(a)] Suppose $k$ is not of the form $6n$ with $n$ a power of 2. Then $v_2(W_{k,\mu(k)})=\lambda(k), $ and if \(b<\mu(k)\) then $v_2(W_{k,b})>\lambda(k).$
		\item[(b)] Suppose $k=6n$, where $n$ is a power of 2. Then 
			$v_2(W_{k,n})=0$, $v_2(W_{k,0})=1,$ and for $b\notin\{0,n\}$ we have $v_2(W_{k,b})>1$. 
	\end{enumerate}
\end{proposition}

\begin{proof}

We argue by induction on $k$.

(a)  After pairing symmetric terms, there is a distinguished
summand of the form $ dG_PG_Q,$ where $P$ is a power of 2 $Q=k-P$, and $ v_2(k)<v_2(P).$ Therefore $v_2(Q)=v_2(k)$ and $\mu(Q)=\mu(k)$. Moreover, the scalar $d$ satisfies $ v_2(d)+\lambda(Q)=\lambda(k). $

For every other summand  $d'G_AG_B$, with $A+B=k$, the lower-bound analysis gives either
\[
v_2(d')+\lambda(A)+\lambda(B)>\lambda(k),
\]
or else equality holds and $\mu(A)+\mu(B)>\mu(k).$

Expanding
\[ dG_PG_Q=d \left( \sum_u W_{P,u} G_4^{(P-6u)/4} G_6^u \right) \left( \sum_v W_{Q,v} G_4^{(Q-6v)/4} G_6^v \right),\] 
we see that its contribution to \(W_{k,\mu(k)}\) is $d\sum_{u+v=\mu(k)}W_{P,u}W_{Q,v}. $ The term with $u=0$ and $v=\mu(k)=\mu(Q)$ is $dW_{P,0}W_{Q,\mu(Q)}$. Since $P$ is a power of 2, $W_{P,0}$ is a 2-adic unit. By the induction hypothesis, $v_2(W_{Q,\mu(Q)})=\lambda(Q).$ Therefore this term has valuation $v_2(d)+v_2(W_{P,0})+v_2(W_{Q,\mu(Q)}) = v_2(d)+\lambda(Q) = \lambda(k). $ Every other way to obtain total $G_6$-degree $\mu(k)$ has $u>0$ and $v=\mu(k)-u<\mu(k)=\mu(Q).$ Since $P$ is a power of 2, $W_{P,u}$ has positive valuation for $u>0$.  Also, by induction, $v_2(W_{Q,v})>\lambda(Q)$ whenever $v<\mu(Q)$.  Hence every term with  $u>0$ contributes with valuation strictly larger than $v_2(d)+ \lambda(Q)=\lambda(k).$ Thus the total contribution of the distinguished summand to $W_{k,\mu(k)}$ has valuation exactly $\lambda(k)$.

Now, consider a non-distinguished summand $d'G_AG_B$, with $A+B=k$. Its contribution to $W_{k,\mu(k)}$ is $d'\sum_{u+v=\mu(k)}W_{A,u}W_{B,v}$. If
$ v_2(d')+\lambda(A)+\lambda(B)>\lambda(k),$ then every individual term $d'W_{A,u}W_{B,v}$ has valuation \(>\lambda(k)\), because $v_2(W_{A,u})\ge \lambda(A)$, and $v_2(W_{B,v})\ge \lambda(B).$ Therefore this whole summand contributes to \(W_{k,\mu(k)}\) with valuation strictly larger than  $\lambda(k)$.

The remaining possibility is $v_2(d')+\lambda(A)+\lambda(B)=\lambda(k).$
By the equality analysis from the preceding sections: $\mu(A)+\mu(B)>\mu(k).$ For every pair $u,v$ with $u+v=\mu(k),$ we must have either
$ u<\mu(A)$  or $v<\mu(B).$ By the induction hypothesis, this implies either
$v_2(W_{A,u})>\lambda(A)$ or $v_2(W_{B,v})>\lambda(B).$ Thus
$ v_2(W_{A,u}W_{B,v})>\lambda(A)+\lambda(B),$
and hence $v_2(d'W_{A,u}W_{B,v})>\lambda(k).$ So this non-distinguished summand also contributes to $W_{k,\mu(k)}$ with valuation strictly larger than $\lambda(k)$. Therefore the only contribution to $W_{k,\mu(k)}$ with valuation exactly $\lambda(k)$ comes from the distinguished summand, and so $v_2(W_{k,\mu(k)})=\lambda(k).$

The same argument proves the second assertion.  Indeed, if $b<\mu(k)$, then in the distinguished summand every decomposition $u+v=b$ either has $u>0$, or has $u=0$ and $v=b<\mu(Q)$.  In both cases the contribution has valuation strictly larger than $\lambda(k)$.  For every non-distinguished summand, the same argument as above applies, since $ b<\mu(k)<\mu(A)+\mu(B) $ in the equality case.  Hence $v_2(W_{k,b})>\lambda(k) $ for all \(b<\mu(k)\).

\bigskip

(b) The case $k=6$ is immediate, so assume $n\ge 2$. We use the notation from the $k=6n$ identity. This time, the distinguished summand is the middle term $ dG_{3n}^2$, with $d=a_{n/2}+b_{n/2}$, and $v_2(d)=2s(n)-s(3n)=0.$ Also, note that $\mu(3n)=n/2$ and $\lambda(3n)=0.$

By induction, the coefficient $W_{3n,n/2}$ has valuation 0, while every coefficient $W_{3n,u}$ with $u<n/2$ has positive valuation. The contribution of the middle term to $W_{6n,n}$ contains the term $d\,W_{3n,n/2}^2,$ which has valuation 0. All other decompositions $u+v=n$ have at least one of $u,v$ smaller than $n/2$, and hence contribute with positive valuation.  Thus $v_2(W_{6n,n})=0.$

Next consider $W_{6n,0}$.  The middle term contributes $ d\,W_{3n,0}^2,$ which has valuation at least 2, since $v_2(W_{3n,0})>0$.  The boundary term is $b_nG_{4n}G_{2n},$ and, as computed in the $k=6n$ case, $v_2(b_n)=s(6n)-s(n)=1.$ Both $4n$ and $2n$ are powers of 2, so their pure $G_4$-coefficients are 2-adic units.  Therefore the boundary term contributes to $W_{6n,0}$ with valuation exactly 1. For the remaining paired terms $d_jG_{2n+2j}G_{4n-2j}$, with  $d_j=2a_j+b_j+b_{n-j}$ and $j\neq n/2,n$, the computations in the $k=6n$ section give $v_2(d_j)=1.$ Moreover, the two lower weights cannot both have $\lambda=0$.  Indeed, this would force both $j$ and $n-j$ to be powers of 2; hence $j=n/2$, the excluded middle case.  Therefore every coefficient coming from $G_{2n+2j}G_{4n-2j}$ has positive valuation, and so these paired terms contribute to \(W_{6n,0}\) with valuation at least \(2\). Hence $v_2(W_{6n,0})=1.$

Finally, suppose $b\notin\{0,n\}$. The boundary term contributes with valuation $>1$, since producing positive $G_6$-degree requires a positive $G_6$-degree coefficient from $G_{4n}$ or $G_{2n}$, and all such coefficients have positive valuation. The remaining paired terms again have valuation $1$ and product contribution of positive valuation, hence contribute with valuation $>1$. 

It remains to consider the middle term $dG_{3n}^2$. By induction applied to
$3n=6(n/2)$, we have $v_2(W_{3n,n/2})=0$, $v_2(W_{3n,0})=1,$ and $v_2(W_{3n,u})>1$ for every admissible $u\notin\{0,n/2\}$. Thus, for $0<b<n$, every decomposition $u+v=b$ gives valuation $>1$, except possibly when $b=n/2$. In that case, the two borderline terms occur symmetrically and
combine as $2dW_{3n,0}W_{3n,n/2}$, which has valuation at least $2$. Therefore, the middle term also contributes
with valuation $>1$. Hence $v_2(W_{6n,b})>1$ whenever $0<b<n$.

\end{proof}

\end{document}
